\chapter{Diseño e implementación del modelo a desarrollar}
\label{cap:diseno}

Este capítulo articula el modelo de investigación y su materialización en software. Se presenta primero el \emph{diseño metodológico} —el tipo de investigación y el modelo de simulación numérica que la sostiene, las variables y su matriz de consistencia, los casos de estudio y su criterio de selección, el procedimiento de verificación y validación con sus instrumentos, y los criterios de aceptación con que se contrasta la hipótesis— y, a continuación, el \emph{diseño e implementación} de EduFEM, el software que concreta ese modelo. Por su finalidad, el trabajo es de carácter \emph{propositivo} —la propuesta de un artefacto que resuelve un problema formativo—; por su naturaleza se inscribe en la \emph{investigación aplicada} de tipo \emph{tecnológico}, con enfoque \emph{cuantitativo} en su fase de verificación y validación \autocites[p.~75]{garciacordoba2005tecnologica}[p.~7]{sampieri2018metodologia}: la calidad del artefacto se establece midiendo errores numéricos, tensiones, desplazamientos y tasas de convergencia frente a referencias de exactitud conocida.

\section{Diseño metodológico}
\label{sec:metodologia}

La Introducción presentó una visión general del enfoque de investigación; esta sección lo desarrolla siguiendo la estructura de un modelo de investigación: la conceptualización y el alcance del modelo, la definición de las variables, el procedimiento y la toma de observaciones, la validación de esas observaciones y el contraste de la hipótesis con los resultados.

\subsection{Tipo de investigación y modelo de simulación numérica}
\label{sec:proceso-met}

El modelo de investigación es un \emph{modelo de simulación numérica} —una representación ideal del objeto en la que el investigador abstrae y sistematiza las relaciones que considera esenciales \autocite[p.~21]{alvarez_metodologia}—: el objeto de estudio declarado en la Introducción —el análisis de medios continuos en elasticidad lineal plana por el MEF— se estudia resolviendo problemas de ese dominio con el motor de cálculo desarrollado y contrastando sus resultados con referencias de exactitud conocida; el campo de acción —la enseñanza de ese análisis— recibe el aporte en forma de los atributos del artefacto que se verifican por diseño. No hay unidades experimentales, tratamientos ni aleatorización: el motor es determinista, de modo que la confiabilidad del modelo no descansa en réplicas sino en la reproducibilidad exacta de cada corrida y en la convergencia de sus resultados bajo refinamiento de malla; la incertidumbre numérica se cuantifica con las tasas de convergencia y, cuando la referencia no es exacta, con la extrapolación de Richardson de la propia secuencia de mallas \autocite{roache1998verification}.

El trabajo combinó dos ciclos articulados. El \emph{desarrollo del software} siguió un proceso iterativo e incremental: cada componente del canal de cálculo del MEF (\autoref{cap:marco_teorico}) se implementó y se sometió a prueba de regresión antes de integrarse al conjunto. La \emph{evaluación de la corrección} se realizó mediante un esquema de verificación y validación (\autoref{sec:verificacion-validacion}): verificación por el método de soluciones manufacturadas y validación contra la viga de Timoshenko y la membrana de Cook \autocite{roache1998verification,salari2000mms}.

\subsection{Variables e indicadores}
\label{sec:variables-met}

La \autoref{tab:variables} operacionaliza las variables del estudio, es decir, especifica para cada una las operaciones concretas con que se la mide \autocite[p.~137]{sampieri2018metodologia}. La \emph{variable independiente} reúne la definición del modelo estructural —geometría, material y acciones con sus apoyos— y las decisiones de discretización que el investigador manipula: el tipo de elemento (Q4 o Q9) y la densidad de malla (tamaño característico $h$, o $N\times N$). La \emph{variable dependiente} es la respuesta estructural calculada —tensiones, deformaciones, desplazamientos y reacciones— junto con su \emph{exactitud} frente a las soluciones de referencia, medida por las normas de error $L^2$ y $H^1$ y por los errores relativos en tensión y desplazamiento. Las restantes magnitudes —tipo de análisis (tensión o deformación plana) y orden de cuadratura— se mantienen como \emph{variables controladas}.

\begin{table}[htbp]
\centering
\caption{Operacionalización de las variables de la investigación.}
\label{tab:variables}
\begin{tabular}{@{}p{2.7cm}p{3.2cm}p{5.1cm}p{2.5cm}@{}}
\toprule
\textbf{Variable} & \textbf{Dimensión} & \textbf{Indicador} & \textbf{Instrumento} \\
\midrule
Independiente: definición del modelo & Geometría; material; acciones y apoyos & Dominio y malla; $E$, $\nu$, espesor $t$; cargas y restricciones & Guion de cada caso (malla, material, cargas y apoyos) \\
Independiente: discretización & Tipo de elemento; densidad de malla & Q4/Q9; tamaño característico $h$ (o $N\times N$) y número de GDL & Guion de cada caso (tipo de elemento y $N$) \\
Dependiente: respuesta estructural & Tensiones; deformaciones; desplazamientos; reacciones & $\sigma_x,\sigma_y,\tau_{xy},\sigma_{\mathrm{VM}}$; $\varepsilon_x,\varepsilon_y,\gamma_{xy}$; $u,v$ nodales; reacciones en los apoyos & Solucionador del MEF \\
Dependiente: exactitud de la solución & Convergencia; error en desplazamiento; error en tensión & Tasa de convergencia (orden $p$); normas $L^2$ y $H^1$; error relativo (\%) en $\sigma_x$ y en deflexión frente al analítico y a SAP2000 & Guiones de V\&V \\
Controladas & Tipo de análisis; integración & Tensión o deformación plana; orden de cuadratura & Constantes del guion de cada caso \\
\bottomrule
\end{tabular}
\end{table}

\subsection{Matriz de consistencia}
\label{sec:matriz-consistencia}

El \textbf{problema general} —la ausencia de un entorno educativo libre y de código abierto, en español, que exponga de forma transparente, interactiva y verificable el canal de cálculo completo del MEF en elasticidad plana y genere una memoria de cálculo trazable al modelo del alumno— se atiende mediante el \textbf{objetivo general} de desarrollar EduFEM, y se contrasta con la \textbf{hipótesis de diseño} de que tal herramienta es factible y constituye un apoyo plausible al aprendizaje. La \autoref{tab:consistencia} traza la correspondencia entre cada objetivo específico, la pregunta de investigación que responde, las variables o indicadores asociados, el método o instrumento empleado y la evidencia que lo respalda en el documento.

\begin{table}[htbp]
\centering
\caption[Matriz de consistencia de la investigación.]{Matriz de consistencia: correspondencia entre objetivos específicos, preguntas de investigación, variables, métodos y evidencia.}
\label{tab:consistencia}
{\footnotesize
\setlength{\tabcolsep}{4pt}
\renewcommand{\arraystretch}{1.2}
\begin{tabular}{@{}p{3.0cm}p{2.3cm}p{2.7cm}p{2.7cm}p{2.4cm}@{}}
\toprule
\textbf{Objetivo específico} & \textbf{Pregunta asociada} & \textbf{Variable / indicador} & \textbf{Método / instrumento} & \textbf{Evidencia} \\
\midrule
Fundamentar teóricamente el MEF en elasticidad plana & PI-1 a PI-3 & Canal de cálculo cubierto & Revisión de la literatura clásica del MEF y de la V\&V & \autoref{cap:marco_teorico} \\
Implementar el motor MEF 2D Q4/Q9 verificado & PI-1 & Tipo de elemento; exactitud (normas $L^2$, $H^1$) & Implementación pura en NumPy/SciPy; pruebas de regresión & \S\ref{sec:motor}; \autoref{cap:resultados} (MMS) \\
Diseñar la interfaz pre/proceso/post & PI-2 & Organización de la interfaz; flujo de trabajo & Diseño centrado en un lienzo único compartido & \S\ref{sec:pre-proceso} \\
Desarrollar los módulos educativos & PI-2 & Cobertura del canal de cálculo & Capas interactivas superpuestas a la malla real & \S\ref{sec:educacion}; \autoref{cap:resultados} \\
Verificar y validar el software & PI-1, PI-3 & Exactitud; tasa de convergencia; bloqueo por cortante & MMS, viga de Timoshenko, membrana de Cook & \autoref{cap:resultados} \\
Generar memoria de cálculo e interoperabilidad & --- & Trazabilidad del cálculo & Generación de PDF; importación/exportación DXF y CSV & \S\ref{sec:memoria-io} \\
\bottomrule
\end{tabular}}
\renewcommand{\arraystretch}{1.0}
\end{table}

Las preguntas de investigación referidas en la matriz son las enunciadas en la Introducción: \textbf{PI-1}, sobre la precisión verificable del motor frente a soluciones analíticas y comerciales; \textbf{PI-2}, sobre la organización de la interfaz y de los módulos que expone la cadena de cálculo de forma transparente; y \textbf{PI-3}, sobre la evidencia de fenómenos numéricos como el bloqueo por cortante del Q4 frente a la convergencia del Q9.

El primer objetivo específico —la fundamentación teórica— es transversal: sustenta las tres preguntas de investigación y se materializa en el \autoref{cap:marco_teorico}, de ahí que la matriz le asocie las tres. El sexto objetivo específico —generar la memoria de cálculo y ofrecer la interoperabilidad de datos— es de carácter \emph{instrumental}: no deriva de una pregunta de investigación propia, sino que provee la trazabilidad documental y el intercambio de formatos que apoyan la transparencia perseguida por \textbf{PI-2}; de ahí que la \autoref{tab:consistencia} no le asocie ninguna pregunta.

\subsection{Casos de estudio y criterio de selección}
\label{sec:poblacion}

En el esquema clásico, la población sería el universo de problemas de elasticidad lineal plana —en tensión o deformación plana— resolubles con elementos cuadriláteros isoparamétricos. El conjunto de casos que sigue no es una muestra estadística de ese universo, sino una selección \emph{intencional} por valor probatorio —una muestra dirigida, no probabilística \autocite[p.~201]{sampieri2018metodologia}—, que es la práctica establecida en la verificación y validación de códigos de cálculo \autocite{roache1998verification}. Cada caso aporta una referencia de naturaleza distinta: el método de soluciones manufacturadas provee una solución exacta arbitraria, resuelta en cuatro configuraciones (cuadrado unitario y cuadrilátero distorsionado, en tensión y en deformación plana); la viga de Timoshenko simplemente apoyada, una solución analítica de la teoría de la elasticidad y un modelo de cáscara en SAP2000; y la membrana de Cook, un caso de referencia histórico sensible a la distorsión geométrica y al bloqueo por cortante. Cada uno se resuelve con ambos tipos de elemento (Q4 y Q9) sobre mallas de refinamiento creciente, lo que permite observar el comportamiento de las variables dependientes en función de las independientes.

\subsection{Procedimiento, instrumentos y reproducibilidad}
\label{sec:validacion-datos}

Cada caso sigue el mismo procedimiento. Un guion construye la malla y define material, cargas y apoyos; un validador de salud del modelo —una función pura sin dependencias de la interfaz— verifica la consistencia de los datos de entrada antes de resolver: existencia de elementos, restricciones suficientes para suprimir los movimientos de cuerpo rígido, ausencia de nodos referenciados inexistentes y no degeneración de los elementos (determinante Jacobiano positivo); los errores críticos bloquean el cálculo. En la verificación por soluciones manufacturadas la consistencia de los datos está garantizada por construcción: el término fuente se deduce analíticamente de la solución exacta elegida a priori, de modo que las cargas y las restricciones de Dirichlet corresponden exactamente a ella. El sistema se ensambla y resuelve con las mismas funciones del motor que consumen los módulos educativos y la memoria de cálculo, lo que garantiza coherencia entre lo que el alumno ve explicado y lo que el programa resuelve, y las tensiones se recuperan por la cadena descrita en el \autoref{cap:marco_teorico}.

Los instrumentos de medición son los guiones de verificación y validación del proyecto (\texttt{tests/vv\_mms.py}, \texttt{tests/vv\_timoshenko.py} y \texttt{tests/vv\_cook.py}), que computan las normas de error —integradas con un punto de Gauss más por dirección que la rigidez (\autoref{sec:verificacion-validacion}) para no subestimar artificialmente el error—, los errores relativos frente a las referencias y las tasas de convergencia, y depositan los datos crudos en archivos CSV. Dentro de cada caso solo se manipulan el tipo de elemento y la densidad de malla; el tipo de análisis, el material, la geometría, las cargas y los apoyos, el orden de cuadratura ($2\times2$ para Q4 y $3\times3$ para Q9) y las tolerancias numéricas del motor —el cero numérico y el determinante Jacobiano mínimo admisible, constantes del programa— permanecen fijos.

La reproducibilidad se apoya en el determinismo del motor: ante las mismas entradas produce las mismas salidas, y cada caso se regenera por completo al reejecutar su guion con las dependencias declaradas en el repositorio. A ello se suma una batería de pruebas de regresión que verifica la estabilidad del modelo de datos ante operaciones de edición —el ciclo de transformación $\text{Q4}\rightarrow\text{Q9}\rightarrow\text{Q4}$ sin deriva numérica, la equivalencia de resultados con identificadores de nodo no contiguos y la interoperabilidad de formatos— y que se reporta en el \autoref{cap:resultados}.

\subsection{Criterios de aceptación y contraste de la hipótesis}
\label{sec:validacion-resultados}

Los resultados se validan por contraste con referencias de naturaleza distinta. La verificación por soluciones manufacturadas exige que las normas de error decrezcan, al refinar la malla, con las tasas asintóticas teóricas ($\mathcal{O}(h^{p+1})$ en $L^2$ y $\mathcal{O}(h^{p})$ en $H^1$); la validación exige que los desplazamientos y las tensiones concuerden con la solución analítica de la viga de Timoshenko y con el modelo de SAP2000, y que la membrana de Cook reproduzca el valor de referencia junto con el comportamiento característico de cada elemento frente al bloqueo por cortante. Los criterios de aceptación están fijados a priori en los propios guiones, que los evalúan automáticamente y terminan con error si alguno falla: tasas de convergencia observadas dentro de $\pm0{,}5$ de las teóricas en las cuatro configuraciones del MMS (desplazamiento y campo de tensiones recuperado); flecha central de la viga de Timoshenko dentro del $3\,\%$ de la solución analítica con corrección de cortante; y, en la membrana de Cook, desplazamiento del Q9 con $N=8$ dentro del $1{,}5\,\%$ del valor de referencia y flecha del Q4 inferior a la del Q9 con la misma malla, que es la firma del bloqueo por cortante. El \autoref{cap:resultados} reporta los errores efectivamente medidos en cada magnitud, incluidas las componentes secundarias de tensión que no forman parte del criterio.

La hipótesis de diseño enunciada en la Introducción —la factibilidad de construir un software que exponga el canal de cálculo completo del MEF en elasticidad plana de forma transparente, interactiva, numéricamente verificable y coherente con el flujo de trabajo profesional— se contrasta \emph{por diseño}: se establece comprobando que el artefacto reúne esos atributos. Sus afirmaciones subordinadas de carácter cuantitativo —que el motor alcanza precisión verificable (PI-1) y que la herramienta evidencia el bloqueo por cortante (PI-3)— se contrastan numéricamente con los criterios anteriores; la organización de la interfaz y de los módulos (PI-2) se documenta en la \autoref{sec:diseno-edufem}. En la nomenclatura del diseño clásico de investigación, el \emph{aporte} es EduFEM, el \emph{atributo manipulado} es la transparencia e interactividad con que se presenta el canal de cálculo, frente a la opacidad del software profesional y la abstracción de los textos clásicos, y el \emph{efecto esperado} es la observabilidad y contrastabilidad del procedimiento, como apoyo a la comprensión de los fundamentos del MEF, en correspondencia con las variables del problema científico; estos tres elementos no deben confundirse con las variables del experimento numérico de la \autoref{tab:variables}, que son las que efectivamente se miden. La evaluación empírica del impacto educativo con estudiantes excede el alcance temporal del trabajo y se plantea como línea futura.
